\documentclass[11pt]{article}
\usepackage[a4paper,margin=27mm]{geometry}
\usepackage{amsmath,amssymb,amsthm,hyperref}
\newtheorem{proposition}{Proposition}
\newtheorem{lemma}[proposition]{Lemma}
\title{Why the classical multi-splay potential does not certify the proposed third logarithm}
\author{Research note, 2 October 2026}
\date{}
\begin{document}
\maketitle

This is a hand proof about one legal initialized multi-splay state. It
rules out a particular per-request potential inequality. It is not a
lower bound against multi-splay competitiveness, does not establish a
recurring bad sequence, and does not assert reachability of the initial
state from a separately prescribed canonical initialization.

Write $\Phi=\sum_v\log_2 s(v)$ for the classical potential, where
$s(v)$ counts descendants inside $v$'s current auxiliary splay tree.
Let $\mathcal P$ denote the primary preferred-path partition. For a
perfect primary node BST with $n=2^H-1$ keys, use depth labels
$0,\ldots,H$ and the midpoint secondary reference on those labels.
The standard final switch emits the accessed node's depth, including
when it is a reference leaf; it splays the node and its existing
reference-ancestor predecessor and successor inside the appropriate
auxiliary subtrees.

\begin{proposition}\label{prop:failure}
Fix any $\alpha\ge0$ and any function $G$ of the primary path partition.
For arbitrarily large $H$ there is a legal initialized multi-splay state
and a single access satisfying:
\begin{enumerate}
\item the primary partition does not change;
\item there is exactly one mandatory final switch, no ancestor switch,
and one initialized secondary bit change;
\item the access performs $\Omega(\log H)$ rotations and
$\Delta\Phi=\Omega(\log H)>0$.
\end{enumerate}
Consequently no universal constant $C$ can make the per-request bound
\[
 \operatorname{cost}(t)+\alpha\Delta\Phi+\Delta G(\mathcal P)
 \le C(1+\log\log(H+2))(1+p_t+f_t)
\]
hold for every legal initialized state, where $p_t$ counts either the
ancestor switches alone or all mandatory switches, and $f_t$ is the
fixed-owner secondary write count just described.
\end{proposition}

\paragraph{Construction.}
Choose $k=2s+1\ge3$ and $H=2^k-1$. Initialize all primary preferences
rightward. The root preferred path has $H$ nodes and ends at the
maximum key $x$. Store its nodes in a perfectly balanced auxiliary BST;
all other auxiliary shapes can be arbitrary. This is a legal
path-compatible physical BST. Query $x$. There are no wrong ancestor
preferences. The only operation is the final switch at the reference
leaf $x$: splay $x$ to the root, then splay its reference predecessor
below it. Its reference successor does not exist. There are no root
mark changes, so $\mathcal P$ is unchanged.

Only depth $H-1$ is emitted. The secondary reference has $2^k$ leaves,
and label $H-1=2^k-2$ has address $11\cdots110$. Thus all-right
initialized secondary bits incur exactly one change. The ancestor
switch count is zero and the mandatory-switch count is one.

\paragraph{First splay.}
The right spine of the balanced auxiliary consists of
$v_0,v_1,\ldots,v_{2s}=x$. Splaying $x$ uses exactly $s$ right-right
double steps, hence $2s$ rotations. In one such double step write $a,b,c$
for the left-subtree sizes at the grandparent, parent, and target.
The target is the maximum, so it has no right subtree. Cancellation
of the unchanged total-subtree rank gives
\[
 \Delta\Phi_{\mathrm{double}}
 =\log_2\frac{(a+b+1)(a+b+c+2)}{(b+c+2)(c+1)}.
\]
For this balanced initial tree, at the $j$th bottom-up double step,
\[
 b=2^{2j-1}-1,\qquad a=2b+1,\qquad c=b-1.
\]
The equality for $c$ remains true after earlier doubles: they preserve
the size of the original right subtree below the as-yet untouched
parent. Therefore
\[
 \Delta\Phi_{\max}
 =\sum_{j=1}^s\log_2\left(6+\frac4{2^{2j-1}-1}\right)
 >s\log_2 6.
\]

\paragraph{Predecessor splay.}
After the first splay, $x$'s left subtree has right spine
$v_1,v_3,\ldots,v_{2s-1}$. At its $j$th spine node, indexed from zero,
the left subtree and the remaining right subtree have sizes
\[
 A_j=6\cdot4^{s-j-1}-1,
 \qquad R_j=2(4^{s-j-1}-1),\qquad 0\le j<s.
\]
In particular $A_j\ge R_j$. The predecessor is the maximum of this
subtree, at depth $s-1$, so its splay uses $s-1$ rotations.

For completeness, splaying the maximum in a tree whose right-spine
nodes have left-subtree size at least right-subtree size cannot decrease
$\Phi$. Before a double step its as-yet untouched grandparent has
$a\ge b+c+2$, since the lower rotations preserved its right-subtree
size. The displayed double-step ratio is therefore at least one.
A final single left rotation has potential change
\[
 \log_2\frac{a+b+1}{b+1}\ge0,
\]
where $a$ and $b$ are the parent's and target's left-subtree sizes.
These observations apply to the predecessor splay above. Hence
\[
 \operatorname{rotations}=3s-1,\qquad
 \Delta\Phi\ge\Delta\Phi_{\max}>s\log_2 6.
\]
As $\log H=\Theta(s)$ whereas
$1+\log\log(H+2)=\Theta(1+\log(s+1))$, the proposition follows.

\paragraph{Interpretation and limitations.}
The example isolates a cold access that the classical potential does
not pay by decreasing: that potential increases substantially. A
different potential involving the secondary history could store the
needed credit. In particular this one-access family is absorbed by an
$O(n)$ startup allowance in a global competitive inequality, and it
does not disprove the proposed $O(\log^{(3)}n)$ competitive factor.
Its scope is the classical potential, multiplied by any nonnegative
coefficient, plus an arbitrary function of the unchanged primary
partition. A function that also depends on secondary bits or on the
physical auxiliary shape is outside this obstruction.

\paragraph{Deterministic check.}
The companion \texttt{potential-check.py} constructs only the $H$-node
auxiliary and performs the two literal bottom-up splays. For odd
$3\le k\le15$, it checks the exact rotation counts, the first-splay
formula, and nonnegativity of the second-splay potential change. It
does not construct the exponentially larger primary reference and
performs no random search. The asymptotic statement above is a hand
proof, not an inference from these finite checks and not Lean checked.

\end{document}
